% status: experimental — NOT part of the published results
% Input-only body: the containing document supplies ctex/amsmath/amssymb.
\section{物理 Lambertian 子类的否定结论（M13）}
\label{sec:physical-gamma-extension}

固定正定参数空间，令
\[
 M(S)=A+\sum_{k\in S}W_k,\quad
 d(M,W)=\operatorname{tr}\{M^{-1}-(M+W)^{-1}\},\quad
 \gamma=\min_{S\subseteq T,\ x\notin T,\ W_x\ne0}
 \frac{d(M(S),W_x)}{d(M(T),W_x)}.
\]
约定允许 $S=T$，因此 $\gamma\le1$；全零更新时置 $\gamma=1$。
这里的 $A$ 是基线 Schur 信息，\emph{不是}逐灯观测 Jacobian $A_k$。
采用仓库 \texttt{NominalScene} 的固定法向、加性反照率坐标，
\[
 s_{kp}=(n_p^T\ell_k)_+,
 \qquad B_{\phi,k}(p,:)=\rho_p[s_{kp},\ h_{kp}n_p^Tt_{1k},\ h_{kp}n_p^Tt_{2k}],
 \qquad h_{kp}=\mathbf1\{n_p^T\ell_k>0\}.
\]
$A_k=\operatorname{diag}(\sqrt{w_k}s_k)$，
$B_k=\operatorname{diag}(\sqrt{w_k})B_{\phi,k}$，
$u_k=A_k^TB_k$，$M_{0k}=B_k^TB_k$，
\[
 A=\sum_k A_k^TA_k-\sum_k u_k(M_{0k}+\Lambda_k)^{-1}u_k^T,
 \qquad
 W_k=u_k\big[(M_{0k}+\Lambda_k)^{-1}
 -(M_{0k}+\kappa\Lambda_k)^{-1}\big]u_k^T.
\]
下例取 $\rho=1$，所以在该名义点加性反照率与 log-albedo 的一阶 Jacobian 相同。

\paragraph{显式物理参数族。}
取 $0<q<1$，$\varepsilon=2q/(1+q^2)$，
$c=(1-q^2)/(1+q^2)>0$，故 $c^2+\varepsilon^2=1$。
设 $P=2,L=3,\kappa=10,\rho_1=\rho_2=1$，
\[
 n_1=(3/5,4/5,0),\qquad n_2=(-c,\varepsilon,0),
 \qquad (\ell_0,\ell_+,\ell_-)=(e_x,e_y,e_y).
\]
切基按代码选择 $t_{1k}=\ell_k\times e_z$，
$t_{2k}=\ell_k\times t_{1k}$。权重与 shading 为
\[
 w_0=(25/9,1),\quad w_+=w_-=(25/16,1),\qquad
 s_0=(3/5,0),\quad s_+=s_-=(4/5,\varepsilon).
\]
所有权重严格正。灯 $0$ 对第二像素严格背向，而非位于阴影导数的折点。
对应的物理 Jacobian \emph{不是任意指定}，而是由上述几何唯一给出
\[
 B_{\phi,0}=\begin{bmatrix}3/5&-4/5&0\\0&0&0\end{bmatrix},\qquad
 B_{\phi,+}=B_{\phi,-}
 =\begin{bmatrix}4/5&3/5&0\\\varepsilon&-c&0\end{bmatrix}.
\]
定义
\[
 H=\begin{bmatrix}1&3/4\\\varepsilon&-c\end{bmatrix},\qquad
 S_\pm=\frac1{20}\begin{bmatrix}101&\pm99\\\pm99&101\end{bmatrix},
\]
以及正定先验
\[
 \Lambda_0=\operatorname{diag}(25/9,25/9,1),\qquad
 \Lambda_\pm=\operatorname{diag}(H^TS_\pm^{-1}H,1).
\]
此处最后的 $\operatorname{diag}$ 表示 $2+1$ 分块直和。
$S_\pm$ 的特征值为 $10,1/10$，且
$\det H=-c-3\varepsilon/4\ne0$，故全部先验严格正定。
有理 $q$ 给出全部有理的 $n,\ell,w,s,B_\phi,\Lambda$。
先验允许逐灯异质及 log-intensity--方向相关性；不声称反例已覆盖
共享且对角、方向各向同性的更窄先验族。

\paragraph{LightBlocks 的精确化简。}
灯 $\pm$ 的白化方向矩阵的前两列正是 $H$，
令 $D=\operatorname{diag}(1,\varepsilon)$。利用
\[
 I-H(H^TH+tH^TS_\pm^{-1}H)^{-1}H^T
 =t(tI+S_\pm)^{-1}
\]
可得该灯的 Schur 贡献为 $D\,t(tI+S_\pm)^{-1}D$。
两灯在 $t=1$ 的和恰为 $D^2$：$S_+$ 与 $S_-$ 交换
$10,1/10$ 的两个特征方向，而 $1/11+10/11=1$。
灯 $0$ 的有效白化行是 $(1,-4/3)$，其先验满足
$b_0^T\Lambda_0^{-1}b_0=1$，故该灯贡献
$\operatorname{diag}(t/(1+t),0)$。
于是实际 \texttt{LightBlocks} 装配给出
\begin{align}
 A&=\operatorname{diag}(3/2,\varepsilon^2),
 &W_0&=\operatorname{diag}(9/22,0),\\
 W_\pm&=\begin{bmatrix}v&\pm z\varepsilon\\
                    \pm z\varepsilon&v\varepsilon^2\end{bmatrix},
 &v&=\frac{99}{404},\quad z=\frac{729}{4444}.
\end{align}
因为 $g(s)=9s/[(1+s)(10+s)]$ 满足
$g(10)=9/22$、$g(1/10)=90/1111$，故
$v=[g(10)+g(1/10)]/2$、$z=[g(10)-g(1/10)]/2$。
所有 $\varepsilon>0$ 均有 $A\succ0$，不使用奇异端点或伪逆裁维。
此外
\[
 A^{-1/2}(W_0+W_++W_-)A^{-1/2}
 =\operatorname{diag}(666/1111,99/202)\prec9I,
\]
故已知共享模型约束 $\sum W_k\preceq(\kappa-1)A$ 也被严格满足。

\paragraph{命题 M13：固定 $\kappa$ 与固定 $\alpha$ 仍可 $\gamma\to0$。}
该物理族满足
\[
 \alpha=\max_k\lambda_{\max}(A^{-1/2}W_kA^{-1/2})
 =\frac{165}{808}+\frac{3\sqrt{172105}}{8888},
\]
与 $\varepsilon$ 无关；灯 $0$ 的白化更新特征值为 $3/11,0$。
独立的有理反演给出
\begin{align}
 d(A,W_0)&=\frac17,\\
 \gamma&\le R(\varepsilon)
 :=\frac{d(A,W_0)}{d(A+W_+,W_0)}
 =\frac{3025408256\,\varepsilon^2}
 {77(30614089\,\varepsilon^2+531441)}\\
 &\le\frac{3025408256}{40920957}\,\varepsilon^2
 \longrightarrow0.
\end{align}
事实上 $R(\varepsilon)=(3025408256/40920957)\varepsilon^2+O(\varepsilon^4)$。
故不存在仅依赖本族固定 $\alpha$、固定 $\kappa=10$ 的正统一下界；
基线条件数 $3/(2\varepsilon^2)$ 的恶化不能忽略。
这是物理子类的否定结论，而非仅一般线性模型的实现。

取 $q=1/100$、$\varepsilon=200/10001$，精确穷举全部 $27$ 个
$(S,T,x)$ 三元组（包含 $S=T$）得到
\[
 \gamma=\frac{121016330240000}{4187205554180957}
 =0.0289014543647527\ldots,
 \qquad \alpha=0.344235619878634\ldots .
\]
最小值由 $x=0,S=\varnothing,T=\{+\}$ 或 $\{-\}$ 达到；此时
\[
 d(A+W_+,W_0)=\frac{598172222025851}{121016330240000}.
\]
无需小数精度也可判否：$\gamma<3/100<1/2$，而 $\alpha<1$，
从而 $1/(1+\alpha)>1/2$。

\paragraph{物理可见性与秩。}
$xy$ 平面只是坐标选择。向量
$v_{\rm cam}=(\varepsilon/(2c),1,\delta)$（$\delta>0$）
与两法向、两灯方向的点积均严格正，特别是
$n_2^Tv_{\rm cam}=\varepsilon/2$。
共同刚体旋转使其归一化方向成为相机 $z$ 轴，即得到
$n_z>0,\ell_z>0$ 的共同开放半球；同步旋转切基保持全部点积和信息块。
两灯位置重复可视为两个独立标定的光源/曝光，模型不要求方向互异；
本有限反例的严格不等式也由连续性保持于足够小的合法方向扰动。
第三 nuisance 坐标有有限正精度但零 Jacobian，是合法的秩亏观测。
耦合矩阵实际秩为 $2$，因此保留率恰等于 $1$ 的模式数为 $P-2=0$，
不能把一般下界 $\max(P-3L,0)$ 当成必须取等号的维数约束。

\paragraph{为什么旧一般反例不能直接搬用；一条已排除路线。}
一般线性模型使用跨像素混合的观测行及任意 nuisance-free 补块，
其观测 Jacobian 不是每灯对角，不能直接宣称物理可实现。
更具体地，若仅保留标量 log-intensity nuisance，
$s\ge0,\rho,w>0$ 强制 $u_{kp}=w_{kp}\rho_ps_{kp}^2\ge0$。
此时基线非对角元为
$A_{pq}=-\sum_k u_{kp}u_{kq}/(M_{0k}+\lambda_k)$。
若要求 $A$ 对角，则每个非负乘积必须分别为零，
从而所有更新也对角、$\gamma=1$；不可能得到非零混合更新。
因此“只用强度、直接保持对角基线”的精确搬运路线已被排除。
本构造通过真实切向导数及相关先验的相反交叉项规避此阻塞，
不是任意解耦 $s$ 与 $B_\phi$。

\section{更强的全迹谱证书及限定意义下的紧度（M15）}
\label{sec:gamma-tightness-extension}

只保留非零候选的索引集 $\mathcal U_+$，记 $m=|\mathcal U_+|$。
零候选既不改变信息也不改变任何边际，可从集合问题完全删去。
对于 $m\ge2$，令 $a_0=\lambda_{\min}(A)$，并定义
\begin{align}
 b_{\rm full}&=\frac{a_0}{\lambda_{\max}(M(\mathcal U_+))},\\
 b_1&=\frac{a_0}{\max_x\lambda_{\max}(M(\mathcal U_+\setminus\{x\}))},\\
 b_2&=\frac{a_0}{\max_{x\ne y}\lambda_{\max}(M(\mathcal U_+\setminus\{x,y\}))}.
\end{align}
$b_1$ 是已经证明的 leave-one 谱界。新的证书为
\begin{equation}
 \boxed{\qquad \gamma\ge
 \max\left\{b_2,\frac{4b_{\rm full}}{(1+b_{\rm full})^2}\right\}
 \ge b_1.\qquad}
 \label{eq:gamma-extension-combined}
\end{equation}
$m=0,1$ 时直接 $\gamma=1$。

\paragraph{leave-two 证明。}
$S=T$ 的比值为 $1$。否则选取非零候选 $y\in T\setminus S$，则
$M(S)\preceq M(\mathcal U_+\setminus\{x,y\})$，而 $M(T)\succeq A$。
由已证成对谱界
$d(M,W)/d(N,W)\ge\lambda_{\min}(N)/\lambda_{\max}(M)$
得到该比值至少为 $b_2$。显然 $b_2\ge b_1$。
这里实际使用了严格包含的离散结构，不是把原界重复记号。

\paragraph{Kantorovich 分量的证明。}
先证一个不依赖逆平方保序的矩阵比较。若
$0\prec X\preceq Y$ 且 $m_0I\preceq X\preceq M_0I$，则
\begin{align}
 X^2&\preceq(m_0+M_0)X-m_0M_0I\\
 &\preceq(m_0+M_0)Y-m_0M_0I
 \preceq\frac{(m_0+M_0)^2}{4m_0M_0}\,Y^2.
\end{align}
首步由 $X$ 的谱分解得到；末步对 $Y$ 的每个特征值使用完全平方
\[
 \frac{(m_0+M_0)^2}{4m_0M_0}
 \left(y-\frac{2m_0M_0}{m_0+M_0}\right)^2\ge0.
\]
给定合法三元组，设 $M=M(S),N=M(T),W=W_x$，对 $0\le t\le1$ 取
$X=(N+tW)^{-1}$、$Y=(M+tW)^{-1}$。
全部中间信息矩阵位于 $A$ 与 $M(\mathcal U_+)$ 之间，故可取
$m_0=1/\lambda_{\max}(M(\mathcal U_+))$、$M_0=1/a_0$。
由导数恒等式
\[
 d(P,W)=\int_0^1\operatorname{tr}\{W(P+tW)^{-2}\}\,dt
\]
及与 PSD 权重 $W$ 取迹、积分，得到
$d(N,W)\le[(1+b_{\rm full})^2/(4b_{\rm full})]d(M,W)$。
整个证明\emph{没有}使用错误的逆平方算子单调性。
因为路径包含 $tW_x$，分母必须使用完整集合；
不把 $b_{\rm full}$ 擅自替换成 $b_1$。

固定共享模型精度倍率 $\kappa$ 与基线条件数
$\chi=\lambda_{\max}(A)/\lambda_{\min}(A)$ 时，
$M(\mathcal U_+)\preceq\kappa A$ 给出进一步的普遍推论
\[
 \gamma\ge\frac{4\kappa\chi}{(\kappa\chi+1)^2}.
\]
此处没有声称该固定 $(\kappa,\chi)$ 常数达到最坏值。
物理 M13 族固定 $\kappa$ 但 $\chi\to\infty$，因此不与该推论矛盾。
对任意 $m$，$W_i\preceq\alpha A$ 也给
$b_2\ge1/[\chi(1+(m-2)\alpha)]$。

\paragraph{可证明的精确最坏值：两候选、固定条件数。}
对于 $m=2$，$b_2=1/\chi$。
固定 $0<a<1$，在二维令 $\chi=1/a$，并取参数 $h>0$
\[
 A=\operatorname{diag}(1,a),\quad
 W_0=h^{-4}\begin{bmatrix}1&0\\0&0\end{bmatrix},\quad
 W_1=h^{-4}\begin{bmatrix}1&h\\h&h^2\end{bmatrix}.
\]
除 $S=T$ 外只有两个比值；有理反演给出
\begin{align}
 R_0(h)&=\frac{(ah^4+a+h^2)(ah^6+2ah^2+h^4+1)}
 {(h^4+1)(a^2h^6+2ah^4+h^2+1)},\\
 R_1(h)&=\frac{(a^2+h^2)(h^4+1)(ah^6+2ah^2+h^4+1)}
 {(ah^4+a+h^2)(a^2h^6+h^8+2h^4+1)}.
\end{align}
因此
\[
 R_0=a+(2a^2-a+1)h^2+O(h^4),\quad
 R_1=a+(2a^2-1+a^{-1})h^2+O(h^4).
\]
两二阶系数之差为 $(a-1)^2/a>0$；固定 $a<1$ 时足够小的 $h$
有 $R_0<R_1<1$，所以
\[
 \gamma=a+(2a^2-a+1)h^2+O(h^4),\qquad
 \inf_{\substack{\dim=2,\ m=2\\ \operatorname{cond}(A)=\chi}}\gamma=\frac1\chi.
\]
该族中 $b_1=ah^4+O(h^6)$，$b_{\rm full}=\Theta(h^4)$，
式 \eqref{eq:gamma-extension-combined} 在足够小 $h$ 时恰返回 $a$。
因此新证书的相对紧度为
$\gamma/a=1+(2a-1+a^{-1})h^2+O(h^4)$，而旧 leave-one 界相对松弛发散。
$a=1$ 端点应单独处理：$b_2=1$，因 $S=T$，$\gamma\equiv1$；
不能把上面的非恒定渐近展开继续用于该端点。
本族的更新范数为 $\Theta(h^{-4})$，故它不解决同时固定
$\kappa$、更新范数、候选数及条件数时的完整最坏值问题，
也不证明 Kantorovich 分量的成对常数最优。

\paragraph{物理族上的差距及收敛阶。}
M13 族有
$\lambda_{\max}(A+W_0+W_++W_-)=5331/2222$
（对足够小 $\varepsilon$），$\lambda_{\min}(A)=\varepsilon^2$，
故新的 Kantorovich 分量为
$(8888/5331)\varepsilon^2+O(\varepsilon^4)$。
结合 $b_2$ 与 $R(\varepsilon)$ 可得
$\gamma=\Theta(\varepsilon^2)$，并且
\[
 \lim_{\varepsilon\to0}\frac{R(\varepsilon)}{
  \max\{b_2,4b_{\rm full}/(1+b_{\rm full})^2\}}
 =\frac{672018808864}{15154394409}=44.3448144958\ldots .
\]
这是该物理族的\emph{阶紧度与上界差距}，不是声称常数 $44.34$ 最优，
也不依赖随机搜索证明紧度。

\paragraph{低秩谱包围与真实设计对照。}
为避免对每个移除集合做 $P\times P$ 特征分解，
对总矩阵 $E=M(\mathcal U_+)$ 只求一次顶部 $r$ 维特征空间 $Q$。
移除因子 $V$（一灯或两灯）后，对 $E-VV^T$ 的分块记
\[
 a_Q=\lambda_{\max}(Q^T(E-VV^T)Q),\quad
 c_Q=\lambda_{r+1}(E),\quad
 b_Q=\|Q^TVV^TQ_\perp\|_2.
\]
Rayleigh 原理与二块比较给出严格的精确算术包围
\[
 a_Q\le\lambda_{\max}(E-VV^T)\le
 \min\left\{\lambda_{\max}(E),
 \frac{a_Q+c_Q+\sqrt{(a_Q-c_Q)^2+4b_Q^2}}2\right\}.
\]
其中 $Z=Q^TV$ 满足
$b_Q^2=\lambda_{\max}\{Z(V^TV-Z^TZ)Z^T\}$，只需求小 Gram 矩阵。
实现的上下端点加入浮点安全余量；不把普通 LAPACK 运算称为区间算术。
低端谱界才是返回的保证，上端用于报告数值包围宽度。
$P$ 很小的紧族验证改从 $A$ 重装 leave-one/two 矩阵，避免大更新相减消去基线。

复用既有加载器、corrected 噪声约定及
\texttt{default\_rng([20260910,object\_index])}，对固定 $11$ 对象各测
level $0.1,0.5$，共 $22$ 个设计点；每场景 $P=1200$、$48$ 个活跃灯，
其余 $94$ 个零更新占位不参与有效候选最大值。
新保证相对已证 leave-one 保证提升至少 $3.9168$ 至 $3.9999$ 倍。
这里对照的是有效谱界，不是已经否定的 $1/(1+\alpha)$ 表达式；
也没有穷举真实场景的集合格或声称这些设计的真实 $\gamma$ 已知。
有理反例、参数族、全部 $22$ 行及随机实现审计
均只写入新的
\texttt{results/theory\_extension\_20260918/gamma\_extensions.json}。
其中随机 $40$ 个小实例、$1839$ 个完整三元组零违反，只用于实现反驳测试，
不作为上述紧度命题的证明。

\paragraph{任务权重与剩余边界。}
所有上述系数都针对未加权全迹及固定正定空间。
若 $Q_H=H^TH\succ0$，可以对
$\widetilde M=Q_H^{-1/2}MQ_H^{-1/2}$ 重新计算证书，因为
$\operatorname{tr}(Q_HM^{-1})=\operatorname{tr}(\widetilde M^{-1})$；
不允许复用原坐标的未加权常数。
秩亏任务可有零分子而正分母，故没有这样的正统一证书。
例如 $M=I_2,N=I_2+\tfrac12\mathbf1\mathbf1^T,
W=\operatorname{diag}(1,0),H=(0,1)$ 给
$d_H(M,W)=0$、$d_H(N,W)=1/28$。
未解问题明确保留：同质对角物理先验的完整判定，
联合固定 $\kappa,m,\chi,\|W\|$ 时的最坏值，及新的成对最优常数。
